\section{Applications of Adaptive Filtering}

\subsection{Introduction}

\textbf{Adaptive filtering} plays a crucial role in modern signal processing applications where the signal characteristics or the environment may \textbf{change over time}.
Unlike fixed filters, adaptive filters are capable of automatically \textbf{adjusting their parameters} based on the input data.
This dynamic nature makes them well suited for scenarios involving time-varying signals, interference suppression or system identification.

\subsection{Noise cancellation}

Let us consider a common scenario in which a microphone recording presents \textbf{noise} in addition to the target \textbf{signal of interest}.
A primary input $d(n)$ is composed of both target signal and interfering noise.
A reference input $u(n)$ captures a version of the noise, which is processed by a \textbf{filter} $W(z)$ to produce the output $y(n)$.
This output is subtracted from the primary input to describe the \textbf{error signal} $e(n) = d(n) - y(n)$, which we want to correspond to the target signal alone.

\begin{figure}[H]
    \centering
    \includegraphics[width=0.55\linewidth]{Noise cancellation scheme.png}
    \caption{Noise cancellation: scenario schematic}
\end{figure}

The goal is to \textbf{cancel the interfering component} from the observed signal.
This setting arises in situations where the available signal is a \textbf{mixture} of a desired source and an undesired interference.
The adaptive filter $W(z)$ is adjusted such that its output $y(n)$ approximates the noise contribution in $d(n)$, thus minimizing $e(n)$.

\subsubsection{Methodology}

The methodology for noise cancellation begins by estimating the interfering signal.
This is achieved by placing a \textbf{microphone} near the noisy source in order to record a reference version of the interference.
The \textbf{recorded interference} is then subtracted from the corrupted signal using an \textbf{adaptive filter} $W(z)$.
Finally, the \textbf{residual} error signal $e(n)$ is computed, which ideally contains an \textbf{estimate of the target signal} with the interference removed.

The working principle relies on the \textbf{orthogonality} property of the adaptive filter: the filter is designed such that the resulting error signal is \textbf{uncorrelated} with the filter input.
Consequently, the error signal $e(n)$ captures the component of the desired response that is orthogonal to the reference signal (i.e. the target signal).

It is important to note that one cannot directly extract the target signal by simply subtracting the reference input from the primary input.
This is due to the fact that the noise signal reaching the primary input is the result of an \textbf{unknown filtering process}: it passes through an unknown impulse response before combining with the target (i.e. environment response).
Therefore, the role of adaptive filtering is to model this unknown response and use it to \textbf{predict the disturbance} that contaminates the primary input based on the reference input.

\subsubsection{Example and LMS implementation}

Let us consider an example: a \textit{Formula 1} driver needs to maintain clear communication with the team, which must be achieved without interference from engine or mechanical noise.
As expected, the acoustic noise generated by the engine and mechanical parts is extremely high: it becomes essential to suppress it.
The proposed solution involves placing a microphone at the noise source to record a reference signal $u(n)$ that captures the unwanted interference, process it through an adaptive filter and compare it against the primary microphone signal $d(n)$, which contains both the target signal and the interference.

The \textbf{Least Mean Squares (LMS)} algorithm provides an effective iterative implementation for this purpose:

\begin{enumerate}
  \item The reference input $u$ is filtered using the current filter coefficients $\underline{\mathrm{w}}$:
  \[
  y(n) = \underline{\mathrm{w}}^T(n) \ \underline{\mathrm{u}}(n)
  \]
  
  \item The target signal $e$ is estimated by computing the error between the primary input and the filter output:
  \[
  e(n) = d(n) - y(n)
  \]
  
  \item The filter coefficients $\underline{\mathrm{w}}$ are then updated to minimize this error:
  \[
  \underline{\mathrm{w}}(n+1) = \underline{\mathrm{w}}(n) + \mu \ \underline{\mathrm{u}}(n) \ e(n)
  \]
\end{enumerate}

Notice that the vector $\underline{\mathrm{w}}$ contains the filter coefficients for $W(z)$, while the vector $\underline{\mathrm{u}}$ contains a set of present and past samples of the reference input $u$.

\subsection{Acoustic echo cancellation}

A similar situation occurs when a microphone captures \textbf{echo} components given by the \textbf{environment} acoustic behaviour, which is contextualized in teleconferencing systems.
In this setup, each participant’s voice is captured by a microphone in one conference room and transmitted to a loudspeaker in the remote room.
The loudspeaker signal can travel back to the local microphone after reflecting off walls and surfaces, resulting in multiple delayed and attenuated copies of the original signal.
These echoes are then transmitted again, creating a \textbf{feedback loop} that degrades communication quality.
The goal is to cancel these echoes at the microphone level, ensuring that only the local speaker's voice is transmitted and the remote signal is not re-captured.

\begin{figure}[H]
    \centering
    \includegraphics[width=0.55\linewidth]{Acoustic echo scheme.png}
    \caption{Acoustic echo cancellation: scenario drawing}
\end{figure}

Acoustic echo cancellation can be achieved by means of adaptive filtering.
The signal played by the loudspeaker $u(n)$ propagates through the acoustic environment, which is modelled by an unknown system $H(z)$.
The microphone captures the resulting echo signal $d(n)$, which may interfere with the local speaker's voice.
We introduce an adaptive filter $W(z)$ to remove the echo: it processes the known loudspeaker signal $u(n)$ to produce an estimate $y(n)$ of the echo component.
This estimate is subtracted (as before) from the microphone signal to produce an error signal $e(n) = d(n) - y(n)$.
The filter $W(z)$ is iteratively adjusted to minimize this error, effectively learning the impulse response of the unknown echo path and cancelling its effect.

\begin{figure}[H]
    \centering
    \includegraphics[width=0.55\linewidth]{Acoustic echo filtering.png}
    \caption{Acoustic echo cancellation: adaptive filtering}
\end{figure}

\subsubsection{Example}

Let us consider an example: two people (Jane and Jim) are speaking simultaneously in a reverberant room.
Jane is equipped with a wearable microphone that directly captures her voice, producing the signal $u(n)$.
An environmental microphone captures the overall acoustic environment, which includes contributions from Jane's reverberated voice, Jim's speech $v(n)$ and ambient noise $b(n)$.
Jane’s voice reaches the environmental microphone after passing through a room impulse response $H(z)$.

\begin{figure}[H]
    \centering
    \includegraphics[width=0.6\linewidth]{Two-people echo.png}
    \caption{Example of acoustic echo problem in a room with response $H(z)$}
\end{figure}

The goal is to isolate and estimate Jim’s speech removing both the echo of Jane’s voice and the background noise from the environmental microphone signal: we employ an adaptive echo cancellation scheme.
Jane's voice also reaches the environmental microphone, after the convolution with the room response $H(z)$.
This reverberated signal is then summed with Jim’s speech $v(n)$ and noise $b(n)$ to form the observed signal $d(n)$.
The adaptive filter $W(z)$ allows to estimate the contribution of Jane’s voice $y(n)$, which is subtracted from the observed signal $d(n)$, yielding the error signal $e(n) = d(n) - y(n)$.
The goal is to adapt $W(z)$ such that $y(n)$ closely matches Jane's voice after the room reverberation $H(z)$, minimizing $e(n)$ and leaving primarily Jim’s speech and noise.
This structure allows for dynamic estimation of the echo path and effective suppression of reverberated interference.

\begin{figure}[H]
    \centering
    \includegraphics[width=0.6\linewidth]{Two-people filtering.png}
    \caption{Adaptive filtering $W(z)$ for acoustic echo cancellation}
\end{figure}

The output of the adaptive filter is, as usual:

\[
y(n) = \underline{\mathrm{w}}^T(n) \ \underline{\mathrm{u}}(n)
\]

We remind that $\underline{\mathrm{w}}$ is the vector of the adaptive filter coefficients, while $\underline{\mathrm{u}}(n)$ contains the past values of the reference signal $u$.
The error signal is the difference between estimated echo and actual received signal:

\[
e(n) = y(n) - d(n) = \underline{\mathrm{w}}^T(n) \ \underline{\mathrm{u}}(n) - \left[ \underline{\mathrm{h}}^T(n) \ \underline{\mathrm{u}}(n) + v(n) + b(n) \right]
\]

We remind that $\underline{\mathrm{h}}$ is the true echo path (i.e. room response), $v$ is Jim's speech and $b$ is background noise.
After the transient phase, the adaptive filter converges to the true echo path:

\[
y(n) \longrightarrow \underline{\mathrm{h}}^T \ \underline{\mathrm{u}}(n) \quad \Rightarrow \quad
e(n) \longrightarrow v(n) + b(n)
\]

This residual signal contains the part of $d$ that cannot be estimated from $u$ effectively isolating the contributions of Jim and the noise.
It is important to note that this process is based solely on the signal $d$ to be cleaned: the algorithm does not receive any prior information about the content to be preserved.

\subsubsection{Echo Return Loss Enhancement (ERLE)}

To evaluate the \textbf{performance} of an echo canceller, a common metric is the \textbf{Echo Return Loss Enhancement (ERLE)}, which is expressed in dB.
The ERLE quantifies the ratio of the \textbf{residual echo energy} after cancellation to the energy before cancellation:

\[
\left. \mathrm{ERLE} \right|_{\mathrm{dB}} = -10 \log_{10} \left( \frac{\sum_n \left[e(n) - v(n)\right]^2}{\sum_n d(n)^2} \right)
\]

As this expression makes clear, ERLE is an integral measure: it may not accurately reflect performance in non-stationary conditions, such as speech with pauses or varying background noise.
For this reason, ERLE is generally computed over successive time intervals.
Instantaneous ERLE is typically estimated by evaluating the average values of $[e(n) - v(n)]^2$ and $d(n)^2$ over analysis windows of duration approximately 50 ms.

\subsection{Background noise removal}

Another situation requires selective removal when a noisy signal is added to the target signal of interest (i.e. hissing noise in analog recordings).
We define the observed signal $x(n) = s(n) + v(n)$ as the sum of the \textbf{desired signal} $s$ and some \textbf{additive noise} $v$.
This noisy input $x$ is processed by an \textbf{adaptive filter} $W(z)$ in order to obtain an estimate $y$ of the original signal $s$.

\begin{figure}[H]
    \centering
    \includegraphics[width=0.6\linewidth]{Background noise removal.png}
    \caption{Background noise removal}
\end{figure}

The model is based on two key assumptions.
First, the signal $s$ and the noise $v$ are \textbf{statistically independent}.
Second, an estimate of the \textbf{noise power spectral density} $P_v(\omega)$ is available.
These conditions are typical in noise reduction applications, particularly in audio signal processing, where stationary or slowly varying noise models allow for spectral estimation and optimal filtering.

\subsubsection{Wiener filtering problem}

The goal is to adjust the adaptive filter $W(z)$ in order to \textbf{minimize the MSE} $J$ between desired clean signal and filter output:

\[
J = E \left\{ \abs{s(n) - y(n)}^2 \right\}
\]

In this context, the signal $s$ represents the desired response of a Wiener filter of input $x$ and output $y$.
According to the \textbf{orthogonality principle}, the MSE reaches its minimum value when the estimation error is orthogonal to the observed data:

\[
E \left\{ \left[ s(n) - y(n) \right] \ x(m) \right\} = 0 \qquad \forall \ n, m
\]

This leads to the Wiener–Hopf equations.
Describing the filter output $y$ by means of convolution between input $x$ and impulse response $w$, we obtain the following \textbf{Wiener-Hopf equations}, which define the optimal filter coefficients $w(k)$ that minimize the MSE (assuming all second-order statistics of the signals are known):

\[
y(n) = \sum_k w(k) \ x(n-k) \quad \Rightarrow \quad
r_{sx}(n - m) = E \left\{ \sum_k w(k) \ x(n - k) \ x(m) \right\} = \sum_k w(k) \ r_x(n - k - m)
\]

We remind that $r_{sx}$ and $r_x$ are, respectively, \textbf{cross-correlation} between $s$ and $x$, and \textbf{auto-correlation} of $x$.

\subsubsection{PSDs and uncorrelation}

In the frequency domain, the cross-correlation $r_{sx}(n)$ corresponds to the \textbf{cross-spectrum} $P_{sx}(\omega)$, which is defined by means of product between Fourier transforms:

\[
r_{sx}(n) \quad \overset{\mathcal{F}}{\longleftrightarrow} \quad P_{sx}(\omega) = S^*(\omega)X(\omega)
\]

Similarly, the autocorrelation $r_x(n)$ corresponds to the \textbf{power spectrum} $P_x(\omega)$:

\[
r_x(n) \quad \overset{\mathcal{F}}{\longleftrightarrow} \quad P_x(\omega) = X^*(\omega)X(\omega)
\]

Using these relationships, we define a relationship in frequency domain and obtain the \textbf{optimal filter} $W(\omega)$:

\[
r_{sx}(n - m) = \sum_k w(k) \ r_x(n - k - m) \quad \overset{\mathcal{F}}{\longleftrightarrow} \quad
P_{sx}(\omega) = W(\omega) \ P_x(\omega) \quad \Rightarrow \quad
W(\omega) = \frac{P_{sx}(\omega)}{P_x(\omega)}
\]

Since the input signal $s$ and the noise $v$ are \textbf{uncorrelated}:

\[
r_{sx}(n) =  E \left\{ s(k) \ x(k+n) \right\} = E \left\{ s(k) \left[s(k+n) + v(k+n)\right] \right\}
\]

\[
= E \left\{ s(k) \ s(k+n) \right\} + \cancel{E \left\{ s(k) \ v(k+n) \right\}} = E\left\{ s(k) \ s(k+n) \right\} = r_s(n)
\]

Therefore, in the frequency domain:
\[
P_{sx}(\omega) = P_s(\omega)
\]

Additionally, for the uncorrelation of $s$ and $v$:
\[
P_x(\omega) = P_s(\omega) + P_v(\omega)
\]

It follows that the optimum filter becomes:

\[
W(\omega) = \frac{P_{sx}(\omega)}{P_x(\omega)} = \frac{P_s(\omega)}{P_s(\omega) + P_v(\omega)}
\]

So, to implement the filter, we need to estimate the power spectra of both signal and noise.
The \textbf{noise PSD estimation} $\hat{P}_v(\omega)$ can be evaluated when the signal is absent (i.e. silence) and allows us to obtain the \textbf{signal PSD estimation} $\hat{P}_s(\omega)$, and consequently the optimal filter $W(\omega)$, as:

\[
\hat{P}_s(\omega) = P_x(\omega) - \hat{P}_v(\omega) \quad \Rightarrow \quad
W(\omega) = \frac{\hat{P}_s(\omega)}{\hat{P}_s(\omega) + \hat{P}_v(\omega)}
\]

The above equation does not account for phase: zero-phase filters are commonly used.

\subsubsection{Spectral attenuation factor}

Let us consider two \textbf{extreme cases}:

\[
\left\{
\begin{array}{cc}
    P_s(\omega) \ll P_v(\omega): & \abs{W(\omega)} \rightarrow 0 \\
    P_s(\omega) \gg P_v(\omega): & \abs{W(\omega)} \rightarrow 1
\end{array}
\right.
\]

In the first case, the noise overcomes the signal: all the filter can do is to attenuate the input signal at those frequencies.
In the second case, the signal overcomes the noise: the filter leaves the signal undistorted.

From the above considerations, we can conclude that the Wiener filter behaves as a \textbf{spectral attenuator}, as its absolute value is between 0 and 1.

\subsection{Active noise control}

\textbf{Active Noise Control (ANC)} refers to situations where acoustic \textbf{anti-noise waves} are generated from electronic circuits in order to interfere and cancel undesired acoustic waves.
A common example is the cancellation of noise in narrow ducts, such as ventilation systems.

\begin{figure}[H]
    \centering
    \includegraphics[width=0.5\linewidth]{Active noise control.png}
    \caption{Block diagram of ANC}
\end{figure}

\clearpage
